\mychapter{Wyniki eksperymentów}
\label{chap:wyniki}

W~rozdziałach \ref{chap:warianty-dijkstry} i~\ref{chap:implementacja} opisano
trzy warianty algorytmu Dijkstry oraz wspólną infrastrukturę implementacyjną,
w~obrębie której zostały one porównane. Niniejszy rozdział prezentuje wyniki
eksperymentów przeprowadzonych zgodnie z~metodyką opisaną w~podrozdziale
\ref{sec:impl:metodyka}: na~trzech grafach pieszych
(Gdańsk~Południe, Gdańsk, Londyn; tab.~\ref{tab:grafy_po}), dla trzech wariantów
algorytmu (bazowy, ograniczony progiem~$t$, kubełkowy) oraz dla dwóch trybów
wykonania (sekwencyjny, wieloprocesowy po~kategoriach usług), z~$100$
powtórzeniami na~każdą kombinację. Próg czasu we wszystkich eksperymentach
wynosi $t=15$~min, prędkość marszu $v=5\,\mathrm{km/h}$, a~szerokość kubełka
$\Delta=1{,}0$~min (uzasadnienie w~podrozdz.~\ref{subsec:warianty:bucket:delta}).

Dwa pierwsze podrozdziały sprawdzają poprawność wyników i~prezentują same
wynikowe izochrony na~mapach badanych obszarów; trzeci zestawia czasy
wykonania w~postaci pełnej tabeli pomiarowej. Kolejne trzy podrozdziały
interpretują te liczby pod~kątem pytań badawczych
\textbf{P1}--\textbf{P3} sformułowanych w~podrozdziale~\ref{sec:intro:cel}:
wpływu ograniczenia progiem~$t$, zysku z~zastąpienia kopca binarnego kolejką
kubełkową oraz przyspieszenia uzyskiwanego przez zrównoleglenie obliczeń
po~kategoriach usług. Ostatni podrozdział analizuje łączny efekt obu modyfikacji
algorytmicznych i~zrównoleglenia. Podsumowanie odpowiedzi
na~pytania \textbf{P1}--\textbf{P3} wraz z~ograniczeniami metodyki
i~kierunkami dalszych badań znajduje się w~rozdz.~\ref{chap:podsumowanie}.

\section{Walidacja poprawności wyników}
\label{sec:wyniki:poprawnosc}

Zgodnie z~procedurą opisaną w~podrozdz.~\ref{sec:impl:metodyka},
przed przystąpieniem do~analizy czasów porównano zwracane zbiory
$C(t)=\bigcap_k C_k(t)$ oraz wszystkie zbiory cząstkowe $C_k(t)$
między trzema wariantami algorytmu i~między trybem sekwencyjnym
a~wieloprocesowym. Porównanie wykonywane jest na~poziomie pełnych,
posortowanych list identyfikatorów wierzchołków, a~nie samych liczności;
gwarantuje to, że dwa zbiory o~tej samej liczności, lecz różnej zawartości,
zostaną zaklasyfikowane jako niezgodne.

\clearpage

Wszystkie zestawienia dały identyczne zbiory. W~szczególności
liczność $|C(t)|$ przyjmuje wartości $3\,104$ dla Gdańska Południe,
$21\,694$ dla Gdańska oraz $208\,610$ dla Londynu, niezależnie od~wariantu
algorytmu i~od~trybu wykonania. Pozwala to przypisać obserwowane różnice
w~czasie wykonania wyłącznie zmianom samego algorytmu lub trybu wywołania,
a~nie różnicom w~tym, co liczone.

\section{Wizualizacja wynikowych zbiorów $C(t)$}
\label{sec:wyniki:mapy}

Aby skonfrontować liczbowy opis zbiorów $C(t)$ z~ich rzeczywistą strukturą
przestrzenną, wynikowe izochrony zostały zwizualizowane na~tle sieci pieszej
każdego z~trzech badanych obszarów. Ponieważ wszystkie trzy warianty algorytmu
zwracają identyczne zbiory $C(t)$ (sekcja~\ref{sec:wyniki:poprawnosc}),
wystarczy zaprezentować mapy dla jednego z~nich.

Kolorowanie odbywa się na~poziomie krawędzi grafu pieszego, nie pojedynczych
wierzchołków. Konkretnie, dla krawędzi $(u,v)\in E$:
\begin{itemize}
  \item \textcolor[HTML]{2ca02c}{kolor zielony} oznacza, że \emph{oba}
        końce krawędzi należą do~zbioru $C(t)$, tj.\ z~obu końców można
        dotrzeć do~wszystkich kategorii usług w~czasie $\le t$;
  \item \textcolor[HTML]{d62728}{kolor czerwony} oznacza, że co~najmniej
        jeden koniec krawędzi nie należy do~$C(t)$.
\end{itemize}
Taki sposób wizualizacji pozwala odczytać wynik jak~mapę uliczną: zielone
ulice spełniają warunek dostępności 15-minutowej do~wszystkich kategorii
usług, a~czerwone go nie spełniają. Wariant z~kolorowaniem krawędziowym, zamiast
punktowego wyróżnienia samych wierzchołków zbioru $C(t)$, został zainspirowany
wizualizacjami z~pracy Lama~\cite{Lam2024}.

Rysunki \ref{fig:wyniki:mapa_gdansk_poludnie}--\ref{fig:wyniki:mapa_london}
przedstawiają mapy wynikowe trzech badanych obszarów. Na~każdej z~nich
udział zielonych krawędzi w~sieci pieszej dobrze odpowiada wartościom
$|C(t)|/|V|$ z~tab.~\ref{tab:wyniki:pokrycie}: $51{,}9\%$ krawędzi
przy $50{,}8\%$ wierzchołków dla Gdańska Południe, $68{,}7\%$ przy $67{,}8\%$
dla Gdańska oraz $75{,}2\%$ przy $77{,}1\%$ dla Londynu. Drobne rozbieżności
(rzędu 1--2 punktów procentowych) wynikają z~tego, że krawędzie zaliczane są
do~„zielonych” dopiero wtedy, gdy oba ich końce należą do~$C(t)$; rozkład
zgodności obu końców na~poziomie pojedynczych krawędzi zależy więc od~spójności
przestrzennej zbioru $C(t)$.

Z~map widać też regularność, której tabele liczbowe same nie oddają.
Na~każdym z~obszarów zielona część sieci pokrywa się z~zaludnioną,
ciasną tkanką uliczną centrów dzielnic, natomiast czerwone fragmenty
koncentrują się na~obrzeżach (Gdańsk Południe: północno-zachodnia
część obszaru i~rejony przy granicy administracyjnej; Gdańsk: peryferyjne
fragmenty na~zachodzie i~południu; Londyn: pas zewnętrzny Greater London).
Jest to konsekwencja sposobu, w~jaki rozłożone są lokalizacje usług:
ich gęstość rośnie wraz z~gęstością zaludnienia, więc obszary peryferyjne,
ekstensywnie zabudowane lub niezabudowane (parki, tereny przemysłowe,
lotniska, magistrale tranzytowe), pozostają poza zasięgiem $t=15$~min marszu
do~co~najmniej jednej kategorii.

\begin{figure}[ht]
      \centering
      \includegraphics[width=0.85\textwidth]{gdansk_poludnie_15mc_modified.png}
      \caption{Mapa zbioru $C(t)$ dla Gdańska Południe ($t=15$~min, $\Delta=1{,}0$~min).
      Zielone odcinki to fragmenty sieci, których oba końce mieszczą się w~zbiorze
      $C(t)$; czerwone to fragmenty z~co~najmniej jednym końcem poza $C(t)$.}
      \label{fig:wyniki:mapa_gdansk_poludnie}
      \end{figure}
      
      \begin{figure}[ht]
      \centering
      \includegraphics[width=0.85\textwidth]{gdansk_15mc_modified.png}
      \caption{Mapa zbioru $C(t)$ dla Gdańska ($t=15$~min, $\Delta=1{,}0$~min).
      Kolorowanie jak~na~rys.~\ref{fig:wyniki:mapa_gdansk_poludnie}.}
      \label{fig:wyniki:mapa_gdansk}
      \end{figure}
      
      \begin{figure}[ht]
      \centering
      \includegraphics[width=0.85\textwidth]{london_15mc_modified.png}
      \caption{Mapa zbioru $C(t)$ dla Londynu ($t=15$~min, $\Delta=1{,}0$~min).
      Kolorowanie jak~na~rys.~\ref{fig:wyniki:mapa_gdansk_poludnie}.}
      \label{fig:wyniki:mapa_london}
\end{figure}

\clearpage

\section{Czasy wykonania}
\label{sec:wyniki:tabela}

Tabela~\ref{tab:wyniki:wszystkie} zestawia średni czas, odchylenie standardowe oraz
minimum i~maksimum czasu pojedynczej iteracji ze $100$ powtórzeń dla każdej
kombinacji (graf, wariant, tryb). Wszystkie wartości wyrażone są w~sekundach.

\begin{table}[H]
\centering
\caption{Czasy wykonania pojedynczego uruchomienia algorytmu ($100$ powtórzeń),
w~sekundach. Tryb \emph{sekw.}\ oznacza wykonanie jednoprocesowe,
tryb \emph{rów.}\ oznacza wykonanie wieloprocesowe po~kategoriach usług
(podrozdz.~\ref{sec:impl:wykonanie}).}
\label{tab:wyniki:wszystkie}
\begin{tabular}{l l l r r r r}
\hline
\textbf{Obszar} & \textbf{Wariant} & \textbf{Tryb}
& \textbf{średnia} & \textbf{odch.\ st.} & \textbf{min} & \textbf{max} \\
\hline
Gd.\ Płd.\ & bazowy           & sekw.\ & $0{,}03208$ & $0{,}00095$ & $0{,}03073$ & $0{,}03896$ \\
Gd.\ Płd.\ & bazowy           & rów.\  & $0{,}01288$ & $0{,}00120$ & $0{,}01074$ & $0{,}01599$ \\
Gd.\ Płd.\ & ogr.\ $t$        & sekw.\ & $0{,}02944$ & $0{,}00113$ & $0{,}02782$ & $0{,}03777$ \\
Gd.\ Płd.\ & ogr.\ $t$        & rów.\  & $0{,}01264$ & $0{,}00116$ & $0{,}00987$ & $0{,}01521$ \\
Gd.\ Płd.\ & kubełkowy        & sekw.\ & $0{,}02947$ & $0{,}00076$ & $0{,}02814$ & $0{,}03195$ \\
Gd.\ Płd.\ & kubełkowy        & rów.\  & $0{,}01266$ & $0{,}00103$ & $0{,}01055$ & $0{,}01528$ \\
\hline
Gdańsk     & bazowy           & sekw.\ & $0{,}25774$ & $0{,}01078$ & $0{,}22244$ & $0{,}33845$ \\
Gdańsk     & bazowy           & rów.\  & $0{,}12886$ & $0{,}00469$ & $0{,}11926$ & $0{,}14447$ \\
Gdańsk     & ogr.\ $t$        & sekw.\ & $0{,}23582$ & $0{,}00312$ & $0{,}23066$ & $0{,}24602$ \\
Gdańsk     & ogr.\ $t$        & rów.\  & $0{,}11894$ & $0{,}00540$ & $0{,}11011$ & $0{,}14435$ \\
Gdańsk     & kubełkowy        & sekw.\ & $0{,}21285$ & $0{,}00725$ & $0{,}20231$ & $0{,}27138$ \\
Gdańsk     & kubełkowy        & rów.\  & $0{,}07603$ & $0{,}00349$ & $0{,}07217$ & $0{,}09545$ \\
\hline
Londyn     & bazowy           & sekw.\ & $3{,}45087$ & $0{,}03503$ & $3{,}41964$ & $3{,}59322$ \\
Londyn     & bazowy           & rów.\  & $1{,}79633$ & $0{,}07112$ & $1{,}73581$ & $2{,}25290$ \\
Londyn     & ogr.\ $t$        & sekw.\ & $3{,}29730$ & $0{,}03461$ & $3{,}25583$ & $3{,}41621$ \\
Londyn     & ogr.\ $t$        & rów.\  & $1{,}73680$ & $0{,}06804$ & $1{,}67651$ & $2{,}27200$ \\
Londyn     & kubełkowy        & sekw.\ & $2{,}59778$ & $0{,}03815$ & $2{,}54717$ & $2{,}81358$ \\
Londyn     & kubełkowy        & rów.\  & $1{,}03228$ & $0{,}05112$ & $0{,}92697$ & $1{,}30876$ \\
\hline
\end{tabular}
\end{table}

Już z~samych wartości średnich widać dwie ogólne tendencje. Po~pierwsze,
czas wykonania rośnie z~rozmiarem grafu o~dwa rzędy wielkości pomiędzy
Gdańskiem Południe a~Londynem, co odpowiada wzrostowi $|V|$ o~rząd $\sim 44$
oraz $|E|$ o~rząd $\sim 30$ (tab.~\ref{tab:grafy_po}); skala wzrostu jest
zgodna z~oczekiwaniami wynikającymi z~liniowo--logarytmicznej złożoności
algorytmu Dijkstry. Po~drugie, względne uporządkowanie wariantów nie jest
stałe; na~małym grafie różnice między wariantami są w~granicach szumu
pomiarowego, podczas gdy na~Gdańsku i~Londynie wariant kubełkowy odrywa
się od~pozostałych, szczególnie w~trybie wieloprocesowym. Szczegółowa
interpretacja tych obserwacji jest przedmiotem kolejnych podrozdziałów.

\clearpage

\section{Wpływ ograniczenia progiem $t$ (P1)}
\label{sec:wyniki:p1}

Tabela~\ref{tab:wyniki:p1} zestawia czasy wariantu bazowego (klasyczny algorytm
Dijkstry bez ograniczenia progiem) z~wariantem ograniczonym progiem~$t$,
obliczone w~trybie sekwencyjnym. Iloraz $T_{\text{bazowy}}/T_{\text{ogr.}}$
mierzy przyspieszenie uzyskane wyłącznie przez wprowadzenie progu, przy
identycznej reprezentacji grafu i~identycznej kolejce priorytetowej.

\begin{table}[H]
\centering
\caption{Porównanie wariantu bazowego z~wariantem ograniczonym progiem~$t$
(tryb sekwencyjny). $T$~oznacza średni czas pojedynczego uruchomienia
w~sekundach.}
\label{tab:wyniki:p1}
\begin{tabular}{l r r r}
\hline
\textbf{Obszar} & $T_{\text{bazowy}}$ & $T_{\text{ogr.}\,t}$
& \textbf{Przyspieszenie} \\
\hline
Gdańsk Południe & $0{,}03208$ & $0{,}02944$ & $1{,}09\times$ \\
Gdańsk          & $0{,}25774$ & $0{,}23582$ & $1{,}09\times$ \\
Londyn          & $3{,}45087$ & $3{,}29730$ & $1{,}05\times$ \\
\hline
\end{tabular}
\end{table}

Przyspieszenie z~tytułu samego ograniczenia progiem~$t$ okazuje się
bardzo niewielkie: $5\%$ na~Londynie oraz $9\%$ na~obu obszarach Gdańska,
przy czym wartości te mieszczą się w~rzędzie wielkości odchyleń
standardowych z~tab.~\ref{tab:wyniki:wszystkie}. Jest to zgodne z~analizą
z~podrozdz.~\ref{subsec:warianty:modified:zlozonosc}: w~problemie miasta
15-minutowego, z~powodu redukcji wieloźródłowej (podrozdz.~\ref{sec:warianty:supersource}),
izochrona kategorii może pokrywać dominującą część grafu, więc samo ograniczenie
progiem odcina stosunkowo niewielką pracę.
Tabela~\ref{tab:wyniki:pokrycie} obrazuje skalę zjawiska w~liczbach: pokazuje
liczność $|C_k(t)|$ największej kategorii, liczność części wspólnej $|C(t)|$
oraz odpowiadające im udziały w~całkowitej liczbie wierzchołków grafu
(po~wstawieniu usług, tab.~\ref{tab:grafy_po}).

\begin{table}[H]
\centering
\caption{Pokrycie grafu izochronami w~problemie miasta 15-minutowego.
$|C_{\max}|$ oznacza liczność największego zbioru $C_k(t)$ wśród ośmiu
kategorii NEXI. Procenty liczone względem $|V|$ po~wstawieniu usług
(tab.~\ref{tab:grafy_po}).}
\label{tab:wyniki:pokrycie}
\begin{tabular}{l r r r r r}
\hline
\textbf{Obszar} & $|V|$ & $|C_{\max}|$ & $|C_{\max}|/|V|$
& $|C(t)|$ & $|C(t)|/|V|$ \\
\hline
Gdańsk Południe & $6\,106$   & $6\,069$   & $99{,}4\%$ & $3\,104$   & $50{,}8\%$ \\
Gdańsk          & $32\,004$  & $30\,323$  & $94{,}7\%$ & $21\,694$  & $67{,}8\%$ \\
Londyn          & $270\,454$ & $265\,859$ & $98{,}3\%$ & $208\,610$ & $77{,}1\%$ \\
\hline
\end{tabular}
\end{table}

\clearpage

Dla typowych liczności $|S_k|$
spotykanych w~zaludnionym mieście (tab.~\ref{tab:uslugi_kategorie}) próg~$t=15$
minut marszu od~\emph{najbliższej} usługi kategorii jest spełniony niemal
wszędzie. Ograniczenie progiem odcina więc tylko nieliczne wierzchołki na~obrzeżach
i~w~niepołączonych spójnych składowych grafu, co tłumaczy obserwowane
przyspieszenie rzędu kilku procent.

\section{Wpływ kolejki kubełkowej (P2)}
\label{sec:wyniki:p2}

Tabela~\ref{tab:wyniki:p2} zestawia czasy wariantu ograniczonego progiem
z~wariantem kubełkowym. Iloraz $T_{\text{ogr.}}/T_{\text{kubełkowy}}$ mierzy
przyspieszenie wynikające wyłącznie z~zamiany kopca binarnego na~tablicę
kubełków FIFO (oba warianty stosują to samo ograniczenie progiem~$t$
i~tę samą reprezentację grafu).

\begin{table}[H]
\centering
\caption{Porównanie wariantu ograniczonego progiem z~wariantem kubełkowym
(tryb sekwencyjny). Ostatnia kolumna podaje przyspieszenie wynikające
z~zastąpienia kopca binarnego kolejką kubełkową.}
\label{tab:wyniki:p2}
\begin{tabular}{l r r r}
\hline
\textbf{Obszar} & $T_{\text{ogr.}\,t}$ & $T_{\text{kubełkowy}}$
& \textbf{Przyspieszenie} \\
\hline
Gdańsk Południe & $0{,}02944$ & $0{,}02947$ & $1{,}00\times$ \\
Gdańsk          & $0{,}23582$ & $0{,}21285$ & $1{,}11\times$ \\
Londyn          & $3{,}29730$ & $2{,}59778$ & $1{,}27\times$ \\
\hline
\end{tabular}
\end{table}

W~odróżnieniu od~P1 przyspieszenie kubełkowe wyraźnie zależy od~skali
zadania. Na~Gdańsku Południe oba warianty osiągają praktycznie identyczny
czas wykonania ($1{,}00\times$); na~Gdańsku zysk wynosi $11\%$, a~na~Londynie
$27\%$.

Obserwowany wzorzec jest zgodny z~analizą z~podrozdz.~\ref{subsec:warianty:bucket:zlozonosc}.
W~kopcu binarnym koszt operacji wstawienia i~ekstrakcji rośnie jak~$\log|V'|$,
gdzie~$|V'|$ to~liczba elementów aktualnie w~kolejce; w~tablicy kubełków
analogiczne operacje mają amortyzowany koszt stały, niezależny od~$|V'|$.
Dla małych grafów (Gdańsk Południe, $|V'|\approx 3\,000$ wierzchołków
w~izochronie) rzeczywista różnica między $\log_2|V'|\approx 11$ a~stałym
kosztem jest niewielka i~zostaje zrównoważona przez wyższe koszty stałe
operacji w~tablicy kubełków. Dla większych grafów ($|V'|\approx 30\,000$
na~Gdańsku, $|V'|\approx 270\,000$ na~Londynie) różnica $\log_2|V'|$
między $\sim 15$ a~$\sim 18$, po~pomnożeniu przez liczbę operacji rzędu~$|E'|$,
przekłada się na~mierzalne przyspieszenie wykonania, zgodnie z~oczekiwaną
zależnością logarytmiczną. Jak pokażą wyniki w~podrozdz.~\ref{sec:wyniki:p3},
różnica wariantów ujawnia się jeszcze wyraźniej w~trybie wieloprocesowym.

\clearpage

\section{Wpływ zrównoleglenia obliczeń (P3)}
\label{sec:wyniki:p3}

Tabela~\ref{tab:wyniki:p3} zestawia przyspieszenia uzyskane przez tryb
wieloprocesowy względem trybu sekwencyjnego, osobno dla każdego z~trzech
wariantów algorytmu. Przyspieszenie liczone jest jako iloraz
$S=T_{\text{sekw.}}/T_{\text{rów.}}$ średnich czasów wykonania.

\begin{table}[H]
\centering
\caption{Przyspieszenie trybu wieloprocesowego względem sekwencyjnego.
$T_{\text{sekw.}}$ i~$T_{\text{rów.}}$: średnie czasy pojedynczego
uruchomienia w~obu trybach (w~sekundach).}
\label{tab:wyniki:p3}
\begin{tabular}{l l r r r}
\hline
\textbf{Obszar} & \textbf{Wariant}
& $T_{\text{sekw.}}$ & $T_{\text{rów.}}$ & $S$ \\
\hline
Gdańsk Południe & bazowy     & $0{,}03208$ & $0{,}01288$ & $2{,}49\times$ \\
Gdańsk Południe & ogr.\ $t$  & $0{,}02944$ & $0{,}01264$ & $2{,}33\times$ \\
Gdańsk Południe & kubełkowy  & $0{,}02947$ & $0{,}01266$ & $2{,}33\times$ \\
\hline
Gdańsk          & bazowy     & $0{,}25774$ & $0{,}12886$ & $2{,}00\times$ \\
Gdańsk          & ogr.\ $t$  & $0{,}23582$ & $0{,}11894$ & $1{,}98\times$ \\
Gdańsk          & kubełkowy  & $0{,}21285$ & $0{,}07603$ & $2{,}80\times$ \\
\hline
Londyn          & bazowy     & $3{,}45087$ & $1{,}79633$ & $1{,}92\times$ \\
Londyn          & ogr.\ $t$  & $3{,}29730$ & $1{,}73680$ & $1{,}90\times$ \\
Londyn          & kubełkowy  & $2{,}59778$ & $1{,}03228$ & $2{,}52\times$ \\
\hline
\end{tabular}
\end{table}

Mimo że obliczenia zostały rozdzielone na~$K=8$ równoległych procesów
(po~jednym na~kategorię usług), obserwowane przyspieszenie nie przekracza
$2{,}80\times$. Dwa czynniki to wyjaśniają:

\begin{itemize}
  \item \textbf{niezbilansowane obciążenie kategorii.}
        Liczność $|S_k|$ usług różni się między kategoriami nawet o~rząd
        wielkości (tab.~\ref{tab:uslugi_kategorie}); kategorie najliczniejsze
        (\emph{sustenance}, \emph{shops}) wymagają większej liczby relaksacji
        krawędzi niż kategorie najmniej liczne (\emph{entertainment}, \emph{parks}).
        Całkowity czas wieloprocesowego uruchomienia jest zdominowany przez
        kategorię najwolniejszą, podczas gdy szybsze procesy zaczynają
        wcześniej oczekiwać bezczynnie,

  \item \textbf{hybrydowa architektura procesora.}
        Procesor użyty w~eksperymentach (Intel Core~i7-1255U,
        tab.~\ref{tab:impl:srodowisko}) łączy dwa rdzenie wydajne
        z~ośmioma rdzeniami energooszczędnymi. Przy ośmiu pracujących
        procesach większość zadań ląduje na~rdzeniach energooszczędnych,
        których wydajność na~wątek jest niższa od~rdzeni wydajnych,
        i~one też wyznaczają czas najwolniejszego procesu.
\end{itemize}

\clearpage

\paragraph{Lepsza skalowalność wariantu kubełkowego w~trybie równoległym.}
Najbardziej uderzającą obserwacją z~tab.~\ref{tab:wyniki:p3} jest fakt,
że wariant kubełkowy skaluje się równolegle wyraźnie lepiej niż oba warianty
z~kopcem binarnym, i~to nie w~granicach szumu pomiarowego. Dla Gdańska
osiąga on $S=2{,}80\times$ wobec $S\approx 2{,}00\times$ dla wariantów z~kopcem;
dla Londynu odpowiednio $2{,}52\times$ wobec $\approx 1{,}91\times$.
W~efekcie różnica wariantów, niewielka w~trybie sekwencyjnym
(podrozdz.~\ref{sec:wyniki:p2}), wzrasta w~trybie wieloprocesowym: na~Gdańsku
wariant kubełkowy jest o~$36\%$ szybszy od~wariantu ograniczonego progiem
($0{,}076$ vs.\ $0{,}119$~s), a~na~Londynie o~$41\%$ ($1{,}03$ vs.\ $1{,}74$~s).

Kubełkowa struktura kolejki priorytetowej najsilniej skraca czas na~najliczniejszych
kategoriach (tych z~największymi $|E'|$, gdzie czynnik logarytmiczny kopca
binarnego daje największy efekt), a~to właśnie kategoria najwolniejsza
decyduje o~czasie sumarycznym puli wieloprocesowej. Przyspieszenie tej
jednej kategorii redukuje więc czas zarówno sekwencyjny, jak i~równoległy,
ale w~wariancie równoległym efekt jest proporcjonalnie większy, gdyż
pozostałe (szybsze) kategorie i~tak nie wydłużają czasu wykonania.

\section{Łączny efekt modyfikacji}
\label{sec:wyniki:laczne}

Tabela~\ref{tab:wyniki:czasy_najlepsze} zestawia łączny efekt obu
omówionych kierunków optymalizacji, biorąc jako linię odniesienia wariant
bazowy w~trybie sekwencyjnym (klasyczny algorytm Dijkstry, bez ograniczenia
progiem, jednoprocesowy) i~zestawiając go z~kombinacją wariantu kubełkowego
oraz zrównoleglenia po~kategoriach.

\begin{table}[H]
\centering
\caption{Łączne przyspieszenie najlepszej kombinacji (wariant kubełkowy,
tryb wieloprocesowy) względem wariantu bazowego w~trybie sekwencyjnym.}
\label{tab:wyniki:czasy_najlepsze}
\begin{tabular}{l r r r}
\hline
\textbf{Obszar} & $T_{\text{baz.\ sekw.}}$ & $T_{\text{kub.\ rów.}}$
& \textbf{Przyspieszenie} \\
\hline
Gdańsk Południe & $0{,}03208$ & $0{,}01266$ & $2{,}53\times$ \\
Gdańsk          & $0{,}25774$ & $0{,}07603$ & $3{,}39\times$ \\
Londyn          & $3{,}45087$ & $1{,}03228$ & $3{,}34\times$ \\
\hline
\end{tabular}
\end{table}

Na~Gdańsku Południe łączne przyspieszenie wynika praktycznie wyłącznie
ze~zrównoleglenia: zysk algorytmiczny dla tak małego grafu jest niemierzalny.
Na~Gdańsku i~Londynie wariant kubełkowy w~trybie wieloprocesowym osiąga
łączne przyspieszenie powyżej $3{,}3\times$ względem wariantu bazowego
sekwencyjnego, przy czym oba czynniki, algorytmiczny i~równoległościowy,
wnoszą wkład.

Dla Londynu przyspieszenie $3{,}34\times$ rozkłada się na~$1{,}33\times$
z~tytułu zamiany wariantu bazowego na~kubełkowy w~trybie sekwencyjnym
($3{,}45087/2{,}59778$, tab.~\ref{tab:wyniki:wszystkie}) oraz $2{,}52\times$
z~tytułu zrównoleglenia wariantu kubełkowego (tab.~\ref{tab:wyniki:p3},
ostatni wiersz); iloczyn $1{,}33\cdot 2{,}52\approx 3{,}35$.

\clearpage

Analogicznie dla Gdańska przyspieszenie $3{,}39\times$ to~$1{,}21\times$
z~tytułu wariantu algorytmicznego ($0{,}25774/0{,}21285$) i~$2{,}80\times$
ze~zrównoleglenia; iloczyn $1{,}21\cdot 2{,}80\approx 3{,}39$. Warto zwrócić
uwagę, że łączne przyspieszenie dla Gdańska jest minimalnie wyższe niż dla
Londynu, co odzwierciedla wyraźniejszy zysk równoległościowy wariantu
kubełkowego na~tym obszarze ($2{,}80\times$ vs.\ $2{,}52\times$ na~Londynie).

Z~praktycznego punktu widzenia oznacza to, że obie ścieżki optymalizacji
są jednocześnie korzystne dla miast średniej wielkości oraz metropolii,
i~ich łączny efekt redukuje czas pojedynczej analizy $15$-minutowej
z~$0{,}258$~s do~$0{,}076$~s w~przypadku Gdańska oraz z~$3{,}45$~s do~$1{,}03$~s
w~przypadku Londynu, przy identycznych zbiorach wynikowych. Podsumowanie
przedstawionych w~niniejszym rozdziale obserwacji oraz odpowiedzi
na~pytania badawcze \textbf{P1}--\textbf{P3} zostały zebrane
w~rozdz.~\ref{chap:podsumowanie}.
